\chapter*{Prefacio}
\addcontentsline{toc}{chapter}{Prefacio}

Existe una gran cantidad de material dedicado a la mecánica orbital, desde
textos académicos hasta guías orientadas a aplicaciones prácticas y
simuladores. Sin embargo, al estudiar estos recursos es frecuente encontrarse
con dos enfoques: por un lado, desarrollos teóricos que presuponen cierta
formación matemática y, por otro, guías prácticas que presentan fórmulas y
procedimientos sin detenerse demasiado en su origen.

El propósito de este documento es adoptar un enfoque intermedio. La intención
no es solamente reunir las fórmulas necesarias para calcular maniobras
orbitales, sino comprender las ideas matemáticas sobre las que se construyen.
Partiendo de herramientas relativamente elementales de álgebra, geometría y
cálculo vectorial, se desarrollarán progresivamente los conceptos necesarios
para describir una órbita y realizar cálculos sobre ella.

No se pretende desarrollar exhaustivamente cada una de las ramas de la
matemática y la física que intervienen en la mecánica orbital. En su lugar, se
introducirán las herramientas necesarias a medida que sean requeridas,
procurando proporcionar los fundamentos suficientes para comprender su
aplicación y dejando un estudio más general de estos temas a la bibliografía
correspondiente.

Siempre que resulte razonable, las fórmulas utilizadas serán derivadas a
partir de conceptos anteriores. El objetivo no es realizar demostraciones por
el mero hecho de hacerlas, sino evitar que los resultados aparezcan como
fórmulas aisladas cuyo origen o significado no resulta evidente. De esta
manera, además de saber qué ecuación utilizar, se busca comprender por qué
funciona, qué representa cada una de sus cantidades y bajo qué supuestos puede
aplicarse.

Finalmente, estos conceptos se utilizarán para resolver problemas habituales
de la navegación orbital: modificar el apoapsis o el periapsis, cambiar el
plano de una órbita, planificar encuentros entre naves, realizar maniobras de
fase y calcular transferencias entre cuerpos celestes.

Aunque los principios desarrollados pertenecen a la mecánica orbital en
general, el documento está pensado principalmente para su aplicación en
simuladores y videojuegos de vuelo espacial como \textit{Juno: New Origins},
\textit{Kerbal Space Program} y \textit{Orbiter}. Estos proporcionan un entorno
especialmente conveniente para experimentar con los conceptos y comprobar de
forma práctica los resultados obtenidos.
